% Options for packages loaded elsewhere
\PassOptionsToPackage{unicode}{hyperref}
\PassOptionsToPackage{hyphens}{url}
%
\documentclass[
]{article}
\usepackage{amsmath,amssymb}
\usepackage{iftex}
\ifPDFTeX
  \usepackage[T1]{fontenc}
  \usepackage[utf8]{inputenc}
  \usepackage{textcomp} % provide euro and other symbols
\else % if luatex or xetex
  \usepackage{unicode-math} % this also loads fontspec
  \defaultfontfeatures{Scale=MatchLowercase}
  \defaultfontfeatures[\rmfamily]{Ligatures=TeX,Scale=1}
\fi
\usepackage{lmodern}
\ifPDFTeX\else
  % xetex/luatex font selection
\fi
% Use upquote if available, for straight quotes in verbatim environments
\IfFileExists{upquote.sty}{\usepackage{upquote}}{}
\IfFileExists{microtype.sty}{% use microtype if available
  \usepackage[]{microtype}
  \UseMicrotypeSet[protrusion]{basicmath} % disable protrusion for tt fonts
}{}
\makeatletter
\@ifundefined{KOMAClassName}{% if non-KOMA class
  \IfFileExists{parskip.sty}{%
    \usepackage{parskip}
  }{% else
    \setlength{\parindent}{0pt}
    \setlength{\parskip}{6pt plus 2pt minus 1pt}}
}{% if KOMA class
  \KOMAoptions{parskip=half}}
\makeatother
\usepackage{xcolor}
\usepackage[margin=1in]{geometry}
\usepackage{color}
\usepackage{fancyvrb}
\newcommand{\VerbBar}{|}
\newcommand{\VERB}{\Verb[commandchars=\\\{\}]}
\DefineVerbatimEnvironment{Highlighting}{Verbatim}{commandchars=\\\{\}}
% Add ',fontsize=\small' for more characters per line
\usepackage{framed}
\definecolor{shadecolor}{RGB}{248,248,248}
\newenvironment{Shaded}{\begin{snugshade}}{\end{snugshade}}
\newcommand{\AlertTok}[1]{\textcolor[rgb]{0.94,0.16,0.16}{#1}}
\newcommand{\AnnotationTok}[1]{\textcolor[rgb]{0.56,0.35,0.01}{\textbf{\textit{#1}}}}
\newcommand{\AttributeTok}[1]{\textcolor[rgb]{0.13,0.29,0.53}{#1}}
\newcommand{\BaseNTok}[1]{\textcolor[rgb]{0.00,0.00,0.81}{#1}}
\newcommand{\BuiltInTok}[1]{#1}
\newcommand{\CharTok}[1]{\textcolor[rgb]{0.31,0.60,0.02}{#1}}
\newcommand{\CommentTok}[1]{\textcolor[rgb]{0.56,0.35,0.01}{\textit{#1}}}
\newcommand{\CommentVarTok}[1]{\textcolor[rgb]{0.56,0.35,0.01}{\textbf{\textit{#1}}}}
\newcommand{\ConstantTok}[1]{\textcolor[rgb]{0.56,0.35,0.01}{#1}}
\newcommand{\ControlFlowTok}[1]{\textcolor[rgb]{0.13,0.29,0.53}{\textbf{#1}}}
\newcommand{\DataTypeTok}[1]{\textcolor[rgb]{0.13,0.29,0.53}{#1}}
\newcommand{\DecValTok}[1]{\textcolor[rgb]{0.00,0.00,0.81}{#1}}
\newcommand{\DocumentationTok}[1]{\textcolor[rgb]{0.56,0.35,0.01}{\textbf{\textit{#1}}}}
\newcommand{\ErrorTok}[1]{\textcolor[rgb]{0.64,0.00,0.00}{\textbf{#1}}}
\newcommand{\ExtensionTok}[1]{#1}
\newcommand{\FloatTok}[1]{\textcolor[rgb]{0.00,0.00,0.81}{#1}}
\newcommand{\FunctionTok}[1]{\textcolor[rgb]{0.13,0.29,0.53}{\textbf{#1}}}
\newcommand{\ImportTok}[1]{#1}
\newcommand{\InformationTok}[1]{\textcolor[rgb]{0.56,0.35,0.01}{\textbf{\textit{#1}}}}
\newcommand{\KeywordTok}[1]{\textcolor[rgb]{0.13,0.29,0.53}{\textbf{#1}}}
\newcommand{\NormalTok}[1]{#1}
\newcommand{\OperatorTok}[1]{\textcolor[rgb]{0.81,0.36,0.00}{\textbf{#1}}}
\newcommand{\OtherTok}[1]{\textcolor[rgb]{0.56,0.35,0.01}{#1}}
\newcommand{\PreprocessorTok}[1]{\textcolor[rgb]{0.56,0.35,0.01}{\textit{#1}}}
\newcommand{\RegionMarkerTok}[1]{#1}
\newcommand{\SpecialCharTok}[1]{\textcolor[rgb]{0.81,0.36,0.00}{\textbf{#1}}}
\newcommand{\SpecialStringTok}[1]{\textcolor[rgb]{0.31,0.60,0.02}{#1}}
\newcommand{\StringTok}[1]{\textcolor[rgb]{0.31,0.60,0.02}{#1}}
\newcommand{\VariableTok}[1]{\textcolor[rgb]{0.00,0.00,0.00}{#1}}
\newcommand{\VerbatimStringTok}[1]{\textcolor[rgb]{0.31,0.60,0.02}{#1}}
\newcommand{\WarningTok}[1]{\textcolor[rgb]{0.56,0.35,0.01}{\textbf{\textit{#1}}}}
\usepackage{graphicx}
\makeatletter
\def\maxwidth{\ifdim\Gin@nat@width>\linewidth\linewidth\else\Gin@nat@width\fi}
\def\maxheight{\ifdim\Gin@nat@height>\textheight\textheight\else\Gin@nat@height\fi}
\makeatother
% Scale images if necessary, so that they will not overflow the page
% margins by default, and it is still possible to overwrite the defaults
% using explicit options in \includegraphics[width, height, ...]{}
\setkeys{Gin}{width=\maxwidth,height=\maxheight,keepaspectratio}
% Set default figure placement to htbp
\makeatletter
\def\fps@figure{htbp}
\makeatother
\setlength{\emergencystretch}{3em} % prevent overfull lines
\providecommand{\tightlist}{%
  \setlength{\itemsep}{0pt}\setlength{\parskip}{0pt}}
\setcounter{secnumdepth}{-\maxdimen} % remove section numbering
\ifLuaTeX
  \usepackage{selnolig}  % disable illegal ligatures
\fi
\usepackage{bookmark}
\IfFileExists{xurl.sty}{\usepackage{xurl}}{} % add URL line breaks if available
\urlstyle{same}
\hypersetup{
  pdftitle={Final Project Data Processing},
  hidelinks,
  pdfcreator={LaTeX via pandoc}}

\title{Final Project Data Processing}
\author{}
\date{\vspace{-2.5em}}

\begin{document}
\maketitle

Load Library and reading the cleaned census data, where we only included
6 columns: AGEGRP: age (in groups); LFACT: labor force detail; AGEIMM:
age at immigration; Citizen: Citizenship; YRIM: year of Immigration;
IMMSTAT: immigrant status. These are the main columns that we will be
performing our analysis on and try to explore the causal relatioinship.

\begin{Shaded}
\begin{Highlighting}[]
\FunctionTok{rm}\NormalTok{(}\AttributeTok{list=}\FunctionTok{ls}\NormalTok{())}

\FunctionTok{library}\NormalTok{(tidyverse)}
\end{Highlighting}
\end{Shaded}

\begin{verbatim}
## -- Attaching core tidyverse packages ------------------------ tidyverse 2.0.0 --
## v dplyr     1.1.4     v readr     2.1.5
## v forcats   1.0.0     v stringr   1.5.1
## v ggplot2   3.5.0     v tibble    3.2.1
## v lubridate 1.9.3     v tidyr     1.3.1
## v purrr     1.0.2     
## -- Conflicts ------------------------------------------ tidyverse_conflicts() --
## x dplyr::filter() masks stats::filter()
## x dplyr::lag()    masks stats::lag()
## i Use the conflicted package (<http://conflicted.r-lib.org/>) to force all conflicts to become errors
\end{verbatim}

\begin{Shaded}
\begin{Highlighting}[]
\FunctionTok{library}\NormalTok{(haven)}
\FunctionTok{library}\NormalTok{(stargazer)}
\end{Highlighting}
\end{Shaded}

\begin{verbatim}
## 
## Please cite as: 
## 
##  Hlavac, Marek (2022). stargazer: Well-Formatted Regression and Summary Statistics Tables.
##  R package version 5.2.3. https://CRAN.R-project.org/package=stargazer
\end{verbatim}

\begin{Shaded}
\begin{Highlighting}[]
\FunctionTok{library}\NormalTok{(ggplot2)}
\FunctionTok{library}\NormalTok{(dplyr)}

\NormalTok{df }\OtherTok{\textless{}{-}} \FunctionTok{read\_csv}\NormalTok{(}\StringTok{"data/cen\_2021\_clean.csv"}\NormalTok{)}
\end{Highlighting}
\end{Shaded}

\begin{verbatim}
## Rows: 980868 Columns: 4
## -- Column specification --------------------------------------------------------
## Delimiter: ","
## dbl (4): AGEGRP, LFACT, YRIM, IMMSTAT
## 
## i Use `spec()` to retrieve the full column specification for this data.
## i Specify the column types or set `show_col_types = FALSE` to quiet this message.
\end{verbatim}

\begin{Shaded}
\begin{Highlighting}[]
\FunctionTok{head}\NormalTok{(df)}
\end{Highlighting}
\end{Shaded}

\begin{verbatim}
## # A tibble: 6 x 4
##   AGEGRP LFACT  YRIM IMMSTAT
##    <dbl> <dbl> <dbl>   <dbl>
## 1     13    88  2008      88
## 2     11     3  2005       2
## 3     13    14  9999       1
## 4     16    13  9999       1
## 5     18    13  9999       1
## 6     16    13  9999       1
\end{verbatim}

Data cleaning Process: 1) coerce all the PUMF code‐vars to numeric in
one go 2) drop the ``never worked'' \& missing codes in LFACT 3) keep
only citizens or immigrants 4) restrict to the age groups you care about
(AGEGRP 13--19) 5) create the Retired outcome variable (Yi)

\begin{itemize}
\tightlist
\item
  individuals who have never worked \& missing values in LFACT (all the
  code bigger than and including 14)
\item
  neither citizens nor immigrants (IMMSTAT here we only need code =
  1\&2)
\item
  those not aged in the interval we are interested in (beyond 4974, so
  agegrp from 13\textasciitilde19)
\end{itemize}

\begin{Shaded}
\begin{Highlighting}[]
\CommentTok{\# df without useless columns (at the moment)}
\NormalTok{df\_clean }\OtherTok{\textless{}{-}}\NormalTok{ df }\SpecialCharTok{\%\textgreater{}\%}
  \FunctionTok{mutate}\NormalTok{(}\FunctionTok{across}\NormalTok{(}
    \FunctionTok{c}\NormalTok{(AGEGRP, LFACT, IMMSTAT, YRIM),}
    \SpecialCharTok{\textasciitilde{}} \FunctionTok{as.numeric}\NormalTok{(.)}
\NormalTok{  )) }\SpecialCharTok{\%\textgreater{}\%}
  
  \FunctionTok{filter}\NormalTok{(}\SpecialCharTok{!}\NormalTok{ LFACT }\SpecialCharTok{\%in\%} \FunctionTok{c}\NormalTok{(}\DecValTok{14}\NormalTok{, }\DecValTok{88}\NormalTok{, }\DecValTok{99}\NormalTok{)) }\SpecialCharTok{\%\textgreater{}\%}
  
  \FunctionTok{filter}\NormalTok{(IMMSTAT }\SpecialCharTok{\%in\%} \FunctionTok{c}\NormalTok{(}\DecValTok{1}\NormalTok{, }\DecValTok{2}\NormalTok{)) }\SpecialCharTok{\%\textgreater{}\%}
  
  \FunctionTok{filter}\NormalTok{(AGEGRP }\SpecialCharTok{\textgreater{}=} \DecValTok{13}\NormalTok{, AGEGRP }\SpecialCharTok{\textless{}=} \DecValTok{19}\NormalTok{) }\SpecialCharTok{\%\textgreater{}\%}
  
  \FunctionTok{mutate}\NormalTok{(}
    \AttributeTok{Retired =} \FunctionTok{if\_else}\NormalTok{(LFACT }\SpecialCharTok{\%in\%} \FunctionTok{c}\NormalTok{(}\DecValTok{11}\NormalTok{, }\DecValTok{12}\NormalTok{, }\DecValTok{13}\NormalTok{), }\DecValTok{1}\DataTypeTok{L}\NormalTok{, }\DecValTok{0}\DataTypeTok{L}\NormalTok{)}
\NormalTok{  )}
\FunctionTok{head}\NormalTok{(df\_clean)}
\end{Highlighting}
\end{Shaded}

\begin{verbatim}
## # A tibble: 6 x 5
##   AGEGRP LFACT  YRIM IMMSTAT Retired
##    <dbl> <dbl> <dbl>   <dbl>   <int>
## 1     16    13  9999       1       1
## 2     18    13  9999       1       1
## 3     16    13  9999       1       1
## 4     16    13     9       2       1
## 5     16    13  9999       1       1
## 6     13     1  9999       1       0
\end{verbatim}

\section{Quick check:}\label{quick-check}

\begin{Shaded}
\begin{Highlighting}[]
\NormalTok{df\_clean }\SpecialCharTok{\%\textgreater{}\%} 
  \FunctionTok{count}\NormalTok{(LFACT)     }\CommentTok{\# should only see 1–13}
\end{Highlighting}
\end{Shaded}

\begin{verbatim}
## # A tibble: 13 x 2
##    LFACT      n
##    <dbl>  <int>
##  1     1 192662
##  2     2   9486
##  3     3   4816
##  4     4   3782
##  5     5    929
##  6     6   1316
##  7     7    616
##  8     8    204
##  9     9   6251
## 10    10   2159
## 11    11  10642
## 12    12  11850
## 13    13 131197
\end{verbatim}

\begin{Shaded}
\begin{Highlighting}[]
\NormalTok{df\_clean }\SpecialCharTok{\%\textgreater{}\%} 
  \FunctionTok{count}\NormalTok{(IMMSTAT)   }\CommentTok{\# only 1,2}
\end{Highlighting}
\end{Shaded}

\begin{verbatim}
## # A tibble: 2 x 2
##   IMMSTAT      n
##     <dbl>  <int>
## 1       1 270777
## 2       2 105133
\end{verbatim}

\begin{Shaded}
\begin{Highlighting}[]
\NormalTok{df\_clean }\SpecialCharTok{\%\textgreater{}\%} 
  \FunctionTok{count}\NormalTok{(AGEGRP)    }\CommentTok{\# only 13–19}
\end{Highlighting}
\end{Shaded}

\begin{verbatim}
## # A tibble: 7 x 2
##   AGEGRP     n
##    <dbl> <int>
## 1     13 58613
## 2     14 60607
## 3     15 68034
## 4     16 64749
## 5     17 53613
## 6     18 42635
## 7     19 27659
\end{verbatim}

\begin{Shaded}
\begin{Highlighting}[]
\NormalTok{df\_clean }\SpecialCharTok{\%\textgreater{}\%} 
  \FunctionTok{count}\NormalTok{(Retired)   }\CommentTok{\# retired vs. not}
\end{Highlighting}
\end{Shaded}

\begin{verbatim}
## # A tibble: 2 x 2
##   Retired      n
##     <int>  <int>
## 1       0 222221
## 2       1 153689
\end{verbatim}

Separating treatment/ control group

\begin{Shaded}
\begin{Highlighting}[]
\NormalTok{df\_did }\OtherTok{\textless{}{-}}\NormalTok{ df\_clean }\SpecialCharTok{\%\textgreater{}\%}
  \CommentTok{\# define group membership once and for all}
  \FunctionTok{mutate}\NormalTok{(}
    \AttributeTok{eligible =} \FunctionTok{as.integer}\NormalTok{(}
\NormalTok{      IMMSTAT }\SpecialCharTok{==} \DecValTok{1}                          \CommentTok{\# all citizens}
      \SpecialCharTok{|}\NormalTok{ (IMMSTAT }\SpecialCharTok{==} \DecValTok{2} \SpecialCharTok{\&}\NormalTok{ YRIM }\SpecialCharTok{\textless{}=} \DecValTok{2011}\NormalTok{)      }\CommentTok{\# immigrants ≥10 yrs}
\NormalTok{    ),}
    \AttributeTok{post =} \FunctionTok{as.integer}\NormalTok{(AGEGRP }\SpecialCharTok{\textgreater{}} \DecValTok{16}\NormalTok{),        }\CommentTok{\# age 65+ if AGEGRP 17+}
    \AttributeTok{did  =}\NormalTok{ eligible }\SpecialCharTok{*}\NormalTok{ post                 }\CommentTok{\# the actual treatment}
\NormalTok{  )}
\end{Highlighting}
\end{Shaded}

summary statistics

\begin{Shaded}
\begin{Highlighting}[]
\FunctionTok{summary}\NormalTok{(df\_did)}
\end{Highlighting}
\end{Shaded}

\begin{verbatim}
##      AGEGRP          LFACT             YRIM         IMMSTAT    
##  Min.   :13.00   Min.   : 1.000   Min.   :   1   Min.   :1.00  
##  1st Qu.:14.00   1st Qu.: 1.000   1st Qu.:2011   1st Qu.:1.00  
##  Median :16.00   Median : 1.000   Median :9999   Median :1.00  
##  Mean   :15.62   Mean   : 6.125   Mean   :7453   Mean   :1.28  
##  3rd Qu.:17.00   3rd Qu.:13.000   3rd Qu.:9999   3rd Qu.:2.00  
##  Max.   :19.00   Max.   :13.000   Max.   :9999   Max.   :2.00  
##     Retired          eligible           post             did        
##  Min.   :0.0000   Min.   :0.0000   Min.   :0.0000   Min.   :0.0000  
##  1st Qu.:0.0000   1st Qu.:1.0000   1st Qu.:0.0000   1st Qu.:0.0000  
##  Median :0.0000   Median :1.0000   Median :0.0000   Median :0.0000  
##  Mean   :0.4088   Mean   :0.9731   Mean   :0.3296   Mean   :0.3253  
##  3rd Qu.:1.0000   3rd Qu.:1.0000   3rd Qu.:1.0000   3rd Qu.:1.0000  
##  Max.   :1.0000   Max.   :1.0000   Max.   :1.0000   Max.   :1.0000
\end{verbatim}

\begin{Shaded}
\begin{Highlighting}[]
\CommentTok{\# 2. Summarise retirement rates by age and group}
\NormalTok{df\_trends }\OtherTok{\textless{}{-}}\NormalTok{ df\_did }\SpecialCharTok{\%\textgreater{}\%}
  \FunctionTok{group\_by}\NormalTok{(AGEGRP, eligible) }\SpecialCharTok{\%\textgreater{}\%}
  \FunctionTok{summarise}\NormalTok{(}
    \AttributeTok{prop\_retire =} \FunctionTok{mean}\NormalTok{(Retired),}
    \AttributeTok{.groups =} \StringTok{"drop"}
\NormalTok{  )}
\end{Highlighting}
\end{Shaded}

\begin{Shaded}
\begin{Highlighting}[]
\FunctionTok{ggplot}\NormalTok{(df\_trends, }\FunctionTok{aes}\NormalTok{(}\AttributeTok{x =}\NormalTok{ AGEGRP, }\AttributeTok{y =}\NormalTok{ prop\_retire, }\AttributeTok{color =} \FunctionTok{factor}\NormalTok{(eligible))) }\SpecialCharTok{+}
  \FunctionTok{geom\_line}\NormalTok{(}\AttributeTok{size =} \DecValTok{1}\NormalTok{) }\SpecialCharTok{+}
  \FunctionTok{geom\_point}\NormalTok{(}\AttributeTok{size =} \DecValTok{2}\NormalTok{) }\SpecialCharTok{+}
  \FunctionTok{geom\_vline}\NormalTok{(}\AttributeTok{xintercept =} \FloatTok{16.8}\NormalTok{, }\AttributeTok{linetype =} \StringTok{"dashed"}\NormalTok{, }\AttributeTok{color =} \StringTok{"grey50"}\NormalTok{) }\SpecialCharTok{+}
  \FunctionTok{scale\_color\_manual}\NormalTok{(}
    \AttributeTok{values =} \FunctionTok{c}\NormalTok{(}\StringTok{"blue"}\NormalTok{, }\StringTok{"red"}\NormalTok{),}
    \AttributeTok{labels =} \FunctionTok{c}\NormalTok{(}\StringTok{"Ineligible (control)"}\NormalTok{, }\StringTok{"Eligible (treatment)"}\NormalTok{),}
    \AttributeTok{name   =} \StringTok{"OAS Eligibility"}
\NormalTok{  ) }\SpecialCharTok{+}
  \FunctionTok{scale\_x\_continuous}\NormalTok{(}\AttributeTok{breaks =} \DecValTok{13}\SpecialCharTok{:}\DecValTok{19}\NormalTok{) }\SpecialCharTok{+}
  \FunctionTok{labs}\NormalTok{(}
    \AttributeTok{title    =} \StringTok{"Retirement Rates by Age Group and OAS Eligibility"}\NormalTok{,}
    \AttributeTok{subtitle =} \StringTok{"Full series (13–16 = pre{-}trend; jump at 17=65)"}\NormalTok{,}
    \AttributeTok{x        =} \StringTok{"AGEGRP (13=50–54, …, 17=65–69, …)"}\NormalTok{,}
    \AttributeTok{y        =} \StringTok{"Proportion Retired"}
\NormalTok{  ) }\SpecialCharTok{+}
  \FunctionTok{theme\_minimal}\NormalTok{()}
\end{Highlighting}
\end{Shaded}

\begin{verbatim}
## Warning: Using `size` aesthetic for lines was deprecated in ggplot2 3.4.0.
## i Please use `linewidth` instead.
## This warning is displayed once every 8 hours.
## Call `lifecycle::last_lifecycle_warnings()` to see where this warning was
## generated.
\end{verbatim}

\includegraphics{fp_data_processing_files/figure-latex/unnamed-chunk-7-1.pdf}

immigrant only: compare between treatment immigrants and control
immigrants

\begin{Shaded}
\begin{Highlighting}[]
\NormalTok{df\_imm\_only }\OtherTok{\textless{}{-}}\NormalTok{ df\_clean }\SpecialCharTok{\%\textgreater{}\%}
  \CommentTok{\# get rid of citizen}
  \FunctionTok{filter}\NormalTok{(}
\NormalTok{    IMMSTAT }\SpecialCharTok{==} \DecValTok{2}
\NormalTok{  ) }\SpecialCharTok{\%\textgreater{}\%}
  \CommentTok{\# define group membership once and for all}
  \FunctionTok{mutate}\NormalTok{(}
    \AttributeTok{eligible =} \FunctionTok{as.integer}\NormalTok{(YRIM }\SpecialCharTok{\textless{}=} \DecValTok{2011}     \CommentTok{\# immigrants ≥10 yrs}
\NormalTok{    ),}
    \AttributeTok{post =} \FunctionTok{as.integer}\NormalTok{(AGEGRP }\SpecialCharTok{\textgreater{}} \DecValTok{16}\NormalTok{),        }\CommentTok{\# age 65+ if AGEGRP 17+}
    \AttributeTok{did  =}\NormalTok{ eligible }\SpecialCharTok{*}\NormalTok{ post                 }\CommentTok{\# the actual treatment}
\NormalTok{  )}
\end{Highlighting}
\end{Shaded}

\begin{Shaded}
\begin{Highlighting}[]
\CommentTok{\# 2. Summarise retirement rates by age and group}
\NormalTok{df\_trends\_imm }\OtherTok{\textless{}{-}}\NormalTok{ df\_imm\_only }\SpecialCharTok{\%\textgreater{}\%}
  \FunctionTok{group\_by}\NormalTok{(AGEGRP, eligible) }\SpecialCharTok{\%\textgreater{}\%}
  \FunctionTok{summarise}\NormalTok{(}
    \AttributeTok{prop\_retire =} \FunctionTok{mean}\NormalTok{(Retired),}
    \AttributeTok{.groups =} \StringTok{"drop"}
\NormalTok{  )}
\end{Highlighting}
\end{Shaded}

\begin{Shaded}
\begin{Highlighting}[]
\FunctionTok{ggplot}\NormalTok{(df\_trends\_imm, }\FunctionTok{aes}\NormalTok{(}\AttributeTok{x =}\NormalTok{ AGEGRP, }\AttributeTok{y =}\NormalTok{ prop\_retire, }\AttributeTok{color =} \FunctionTok{factor}\NormalTok{(eligible))) }\SpecialCharTok{+}
  \FunctionTok{geom\_line}\NormalTok{(}\AttributeTok{size =} \DecValTok{1}\NormalTok{) }\SpecialCharTok{+}
  \FunctionTok{geom\_point}\NormalTok{(}\AttributeTok{size =} \DecValTok{2}\NormalTok{) }\SpecialCharTok{+}
  \FunctionTok{geom\_vline}\NormalTok{(}\AttributeTok{xintercept =} \FloatTok{16.8}\NormalTok{, }\AttributeTok{linetype =} \StringTok{"dashed"}\NormalTok{, }\AttributeTok{color =} \StringTok{"grey50"}\NormalTok{) }\SpecialCharTok{+}
  \FunctionTok{scale\_color\_manual}\NormalTok{(}
    \AttributeTok{values =} \FunctionTok{c}\NormalTok{(}\StringTok{"blue"}\NormalTok{, }\StringTok{"red"}\NormalTok{),}
    \AttributeTok{labels =} \FunctionTok{c}\NormalTok{(}\StringTok{"Ineligible (control)"}\NormalTok{, }\StringTok{"Eligible (treatment)"}\NormalTok{),}
    \AttributeTok{name   =} \StringTok{"OAS Eligibility"}
\NormalTok{  ) }\SpecialCharTok{+}
  \FunctionTok{scale\_x\_continuous}\NormalTok{(}\AttributeTok{breaks =} \DecValTok{13}\SpecialCharTok{:}\DecValTok{19}\NormalTok{) }\SpecialCharTok{+}
  \FunctionTok{labs}\NormalTok{(}
    \AttributeTok{title    =} \StringTok{"Retirement Rates by Age Group and OAS Eligibility"}\NormalTok{,}
    \AttributeTok{subtitle =} \StringTok{"Full series (13–16 = pre{-}trend; jump at 17=65)"}\NormalTok{,}
    \AttributeTok{x        =} \StringTok{"AGEGRP (13=50–54, …, 17=65–69, …)"}\NormalTok{,}
    \AttributeTok{y        =} \StringTok{"Proportion Retired"}
\NormalTok{  ) }\SpecialCharTok{+}
  \FunctionTok{theme\_minimal}\NormalTok{()}
\end{Highlighting}
\end{Shaded}

\includegraphics{fp_data_processing_files/figure-latex/unnamed-chunk-10-1.pdf}

\subsection{Regression models}\label{regression-models}

\begin{enumerate}
\def\labelenumi{\arabic{enumi}.}
\tightlist
\item
  baseline (additive) model for citizen-immigrant dataset:
\end{enumerate}

\begin{Shaded}
\begin{Highlighting}[]
\NormalTok{lm\_baseline\_c\_i }\OtherTok{\textless{}{-}} \FunctionTok{lm}\NormalTok{(Retired}\SpecialCharTok{\textasciitilde{}}\NormalTok{ post }\SpecialCharTok{+}\NormalTok{ eligible, }\AttributeTok{data=}\NormalTok{df\_did)}

\FunctionTok{summary}\NormalTok{(lm\_baseline\_c\_i)}
\end{Highlighting}
\end{Shaded}

\begin{verbatim}
## 
## Call:
## lm(formula = Retired ~ post + eligible, data = df_did)
## 
## Residuals:
##     Min      1Q  Median      3Q     Max 
## -0.7907 -0.2230 -0.2230  0.2093  0.8236 
## 
## Coefficients:
##             Estimate Std. Error t value Pr(>|t|)    
## (Intercept) 0.176353   0.004106   42.95   <2e-16 ***
## post        0.567677   0.001434  395.88   <2e-16 ***
## eligible    0.046631   0.004163   11.20   <2e-16 ***
## ---
## Signif. codes:  0 '***' 0.001 '**' 0.01 '*' 0.05 '.' 0.1 ' ' 1
## 
## Residual standard error: 0.4125 on 375907 degrees of freedom
## Multiple R-squared:  0.2959, Adjusted R-squared:  0.2959 
## F-statistic: 7.898e+04 on 2 and 375907 DF,  p-value: < 2.2e-16
\end{verbatim}

2.1 did linear probability model for citizen-immigrant dataset:

\begin{Shaded}
\begin{Highlighting}[]
\NormalTok{lm\_did\_c\_i }\OtherTok{\textless{}{-}} \FunctionTok{lm}\NormalTok{(Retired}\SpecialCharTok{\textasciitilde{}}\NormalTok{ post }\SpecialCharTok{*}\NormalTok{ eligible, }\AttributeTok{data=}\NormalTok{df\_did)}

\FunctionTok{summary}\NormalTok{(lm\_did\_c\_i)}
\end{Highlighting}
\end{Shaded}

\begin{verbatim}
## 
## Call:
## lm(formula = Retired ~ post * eligible, data = df_did)
## 
## Residuals:
##     Min      1Q  Median      3Q     Max 
## -0.7907 -0.2230 -0.2230  0.2093  0.8231 
## 
## Coefficients:
##               Estimate Std. Error t value Pr(>|t|)    
## (Intercept)   0.176892   0.004472  39.552   <2e-16 ***
## post          0.564306   0.011186  50.449   <2e-16 ***
## eligible      0.046073   0.004550  10.126   <2e-16 ***
## post:eligible 0.003428   0.011279   0.304    0.761    
## ---
## Signif. codes:  0 '***' 0.001 '**' 0.01 '*' 0.05 '.' 0.1 ' ' 1
## 
## Residual standard error: 0.4125 on 375906 degrees of freedom
## Multiple R-squared:  0.2959, Adjusted R-squared:  0.2959 
## F-statistic: 5.265e+04 on 3 and 375906 DF,  p-value: < 2.2e-16
\end{verbatim}

2.2 did logistic model for citizen-immigrant dataset:

\begin{Shaded}
\begin{Highlighting}[]
\NormalTok{logm\_did\_c\_i }\OtherTok{\textless{}{-}} \FunctionTok{glm}\NormalTok{(Retired}\SpecialCharTok{\textasciitilde{}}\NormalTok{ post }\SpecialCharTok{*}\NormalTok{ eligible, }\AttributeTok{data=}\NormalTok{df\_did, }\AttributeTok{family=}\NormalTok{binomial)}

\FunctionTok{summary}\NormalTok{(logm\_did\_c\_i)}
\end{Highlighting}
\end{Shaded}

\begin{verbatim}
## 
## Call:
## glm(formula = Retired ~ post * eligible, family = binomial, data = df_did)
## 
## Coefficients:
##               Estimate Std. Error z value Pr(>|z|)    
## (Intercept)   -1.53755    0.02841  -54.12   <2e-16 ***
## post           2.58975    0.06346   40.81   <2e-16 ***
## eligible       0.28908    0.02883   10.03   <2e-16 ***
## post:eligible -0.01214    0.06403   -0.19     0.85    
## ---
## Signif. codes:  0 '***' 0.001 '**' 0.01 '*' 0.05 '.' 0.1 ' ' 1
## 
## (Dispersion parameter for binomial family taken to be 1)
## 
##     Null deviance: 508558  on 375909  degrees of freedom
## Residual deviance: 393681  on 375906  degrees of freedom
## AIC: 393689
## 
## Number of Fisher Scoring iterations: 4
\end{verbatim}

\begin{enumerate}
\def\labelenumi{\arabic{enumi}.}
\setcounter{enumi}{2}
\tightlist
\item
  baseline (additive) model for immigrant-only sample:
\end{enumerate}

\begin{Shaded}
\begin{Highlighting}[]
\NormalTok{lm\_baseline\_i }\OtherTok{\textless{}{-}} \FunctionTok{lm}\NormalTok{(Retired}\SpecialCharTok{\textasciitilde{}}\NormalTok{ post }\SpecialCharTok{+}\NormalTok{ eligible, }\AttributeTok{data=}\NormalTok{df\_imm\_only)}

\FunctionTok{summary}\NormalTok{(lm\_baseline\_i)}
\end{Highlighting}
\end{Shaded}

\begin{verbatim}
## 
## Call:
## lm(formula = Retired ~ post + eligible, data = df_imm_only)
## 
## Residuals:
##     Min      1Q  Median      3Q     Max 
## -0.7698 -0.1864 -0.1864  0.2302  0.8262 
## 
## Coefficients:
##             Estimate Std. Error t value Pr(>|t|)    
## (Intercept) 0.173846   0.003987    43.6   <2e-16 ***
## post        0.583360   0.002664   218.9   <2e-16 ***
## eligible    0.012588   0.004196     3.0   0.0027 ** 
## ---
## Signif. codes:  0 '***' 0.001 '**' 0.01 '*' 0.05 '.' 0.1 ' ' 1
## 
## Residual standard error: 0.399 on 105130 degrees of freedom
## Multiple R-squared:  0.3164, Adjusted R-squared:  0.3164 
## F-statistic: 2.433e+04 on 2 and 105130 DF,  p-value: < 2.2e-16
\end{verbatim}

4.1 did linear probability model for immigrant-only sample:

\begin{Shaded}
\begin{Highlighting}[]
\NormalTok{lm\_did\_i }\OtherTok{\textless{}{-}} \FunctionTok{lm}\NormalTok{(Retired}\SpecialCharTok{\textasciitilde{}}\NormalTok{ post }\SpecialCharTok{*}\NormalTok{ eligible, }\AttributeTok{data=}\NormalTok{df\_imm\_only)}

\FunctionTok{summary}\NormalTok{(lm\_did\_i)}
\end{Highlighting}
\end{Shaded}

\begin{verbatim}
## 
## Call:
## lm(formula = Retired ~ post * eligible, data = df_imm_only)
## 
## Residuals:
##     Min      1Q  Median      3Q     Max 
## -0.7706 -0.1860 -0.1860  0.2294  0.8231 
## 
## Coefficients:
##               Estimate Std. Error t value Pr(>|t|)    
## (Intercept)   0.176892   0.004325  40.896   <2e-16 ***
## post          0.564306   0.010818  52.164   <2e-16 ***
## eligible      0.009134   0.004606   1.983   0.0474 *  
## post:eligible 0.020285   0.011162   1.817   0.0692 .  
## ---
## Signif. codes:  0 '***' 0.001 '**' 0.01 '*' 0.05 '.' 0.1 ' ' 1
## 
## Residual standard error: 0.399 on 105129 degrees of freedom
## Multiple R-squared:  0.3165, Adjusted R-squared:  0.3164 
## F-statistic: 1.622e+04 on 3 and 105129 DF,  p-value: < 2.2e-16
\end{verbatim}

4.2 did logistic model for immigrant-only sample:

\begin{Shaded}
\begin{Highlighting}[]
\NormalTok{logm\_did\_i }\OtherTok{\textless{}{-}} \FunctionTok{glm}\NormalTok{(Retired}\SpecialCharTok{\textasciitilde{}}\NormalTok{ post }\SpecialCharTok{*}\NormalTok{ eligible, }\AttributeTok{data=}\NormalTok{df\_imm\_only, }\AttributeTok{family=}\NormalTok{binomial)}

\FunctionTok{summary}\NormalTok{(logm\_did\_i)}
\end{Highlighting}
\end{Shaded}

\begin{verbatim}
## 
## Call:
## glm(formula = Retired ~ post * eligible, family = binomial, data = df_imm_only)
## 
## Coefficients:
##               Estimate Std. Error z value Pr(>|z|)    
## (Intercept)   -1.53755    0.02841 -54.116   <2e-16 ***
## post           2.58975    0.06346  40.809   <2e-16 ***
## eligible       0.06151    0.03019   2.037   0.0416 *  
## post:eligible  0.09809    0.06566   1.494   0.1352    
## ---
## Signif. codes:  0 '***' 0.001 '**' 0.01 '*' 0.05 '.' 0.1 ' ' 1
## 
## (Dispersion parameter for binomial family taken to be 1)
## 
##     Null deviance: 138452  on 105132  degrees of freedom
## Residual deviance: 104739  on 105129  degrees of freedom
## AIC: 104747
## 
## Number of Fisher Scoring iterations: 4
\end{verbatim}

Citizen--Immigrant Models

\begin{Shaded}
\begin{Highlighting}[]
\FunctionTok{library}\NormalTok{(stargazer)}

\FunctionTok{stargazer}\NormalTok{(}
\NormalTok{  lm\_baseline\_c\_i, lm\_did\_c\_i, logm\_did\_c\_i,}
  \AttributeTok{type =} \StringTok{"html"}\NormalTok{,}
  \AttributeTok{out =} \StringTok{"table\_citizen\_immigrant.html"}\NormalTok{,}
  \AttributeTok{title =} \StringTok{"Citizen–Immigrant Models: Retirement Effects"}\NormalTok{,}
  \AttributeTok{column.labels =} \FunctionTok{c}\NormalTok{(}\StringTok{"Baseline"}\NormalTok{, }\StringTok{"DID (LPM)"}\NormalTok{, }\StringTok{"DID (Logit)"}\NormalTok{),}
  \AttributeTok{covariate.labels =} \FunctionTok{c}\NormalTok{(}\StringTok{"Post"}\NormalTok{, }\StringTok{"Eligible"}\NormalTok{, }\StringTok{"Post × Eligible"}\NormalTok{),}
  \AttributeTok{dep.var.labels =} \StringTok{"Retired"}\NormalTok{,}
  \AttributeTok{omit.stat =} \FunctionTok{c}\NormalTok{(}\StringTok{"f"}\NormalTok{, }\StringTok{"ser"}\NormalTok{),   }\CommentTok{\# 忽略 F统计量 和 残差标准误}
  \AttributeTok{report =} \StringTok{"vc*"}\NormalTok{,              }\CommentTok{\# 报告系数和显著性星号}
  \AttributeTok{digits =} \DecValTok{3}\NormalTok{,                  }\CommentTok{\# 小数点后 3 位}
  \AttributeTok{no.space =} \ConstantTok{TRUE}\NormalTok{,}
  \AttributeTok{align =} \ConstantTok{TRUE}\NormalTok{,}
  \AttributeTok{notes =} \StringTok{"Standard errors in parentheses. *** p\textless{}0.01, ** p\textless{}0.05, * p\textless{}0.1"}
\NormalTok{)}
\end{Highlighting}
\end{Shaded}

\begin{verbatim}
## 
## <table style="text-align:center"><caption><strong>Citizen–Immigrant Models: Retirement Effects</strong></caption>
## <tr><td colspan="4" style="border-bottom: 1px solid black"></td></tr><tr><td style="text-align:left"></td><td colspan="3"><em>Dependent variable:</em></td></tr>
## <tr><td></td><td colspan="3" style="border-bottom: 1px solid black"></td></tr>
## <tr><td style="text-align:left"></td><td colspan="3">Retired</td></tr>
## <tr><td style="text-align:left"></td><td colspan="2"><em>OLS</em></td><td><em>logistic</em></td></tr>
## <tr><td style="text-align:left"></td><td>Baseline</td><td>DID (LPM)</td><td>DID (Logit)</td></tr>
## <tr><td style="text-align:left"></td><td>(1)</td><td>(2)</td><td>(3)</td></tr>
## <tr><td colspan="4" style="border-bottom: 1px solid black"></td></tr><tr><td style="text-align:left">Post</td><td>0.568<sup>***</sup></td><td>0.564<sup>***</sup></td><td>2.590<sup>***</sup></td></tr>
## <tr><td style="text-align:left">Eligible</td><td>0.047<sup>***</sup></td><td>0.046<sup>***</sup></td><td>0.289<sup>***</sup></td></tr>
## <tr><td style="text-align:left">Post × Eligible</td><td></td><td>0.003</td><td>-0.012</td></tr>
## <tr><td style="text-align:left">Constant</td><td>0.176<sup>***</sup></td><td>0.177<sup>***</sup></td><td>-1.538<sup>***</sup></td></tr>
## <tr><td colspan="4" style="border-bottom: 1px solid black"></td></tr><tr><td style="text-align:left">Observations</td><td>375,910</td><td>375,910</td><td>375,910</td></tr>
## <tr><td style="text-align:left">R<sup>2</sup></td><td>0.296</td><td>0.296</td><td></td></tr>
## <tr><td style="text-align:left">Adjusted R<sup>2</sup></td><td>0.296</td><td>0.296</td><td></td></tr>
## <tr><td style="text-align:left">Log Likelihood</td><td></td><td></td><td>-196,840.500</td></tr>
## <tr><td style="text-align:left">Akaike Inf. Crit.</td><td></td><td></td><td>393,688.900</td></tr>
## <tr><td colspan="4" style="border-bottom: 1px solid black"></td></tr><tr><td style="text-align:left"><em>Note:</em></td><td colspan="3" style="text-align:right"><sup>*</sup>p<0.1; <sup>**</sup>p<0.05; <sup>***</sup>p<0.01</td></tr>
## <tr><td style="text-align:left"></td><td colspan="3" style="text-align:right">Standard errors in parentheses. *** p<0.01, ** p<0.05, * p<0.1</td></tr>
## </table>
\end{verbatim}

Immigrant-Only Models

\begin{Shaded}
\begin{Highlighting}[]
\FunctionTok{stargazer}\NormalTok{(}
\NormalTok{  lm\_baseline\_i, lm\_did\_i, logm\_did\_i,}
  \AttributeTok{type =} \StringTok{"html"}\NormalTok{,}
  \AttributeTok{out =} \StringTok{"table\_immigrant\_only.html"}\NormalTok{,}
  \AttributeTok{title =} \StringTok{"Immigrant{-}Only Models: Retirement Effects"}\NormalTok{,}
  \AttributeTok{column.labels =} \FunctionTok{c}\NormalTok{(}\StringTok{"Baseline"}\NormalTok{, }\StringTok{"DID (LPM)"}\NormalTok{, }\StringTok{"DID (Logit)"}\NormalTok{),}
  \AttributeTok{covariate.labels =} \FunctionTok{c}\NormalTok{(}\StringTok{"Post"}\NormalTok{, }\StringTok{"Eligible"}\NormalTok{, }\StringTok{"Post × Eligible"}\NormalTok{),}
  \AttributeTok{dep.var.labels =} \StringTok{"Retired"}\NormalTok{,}
  \AttributeTok{omit.stat =} \FunctionTok{c}\NormalTok{(}\StringTok{"f"}\NormalTok{, }\StringTok{"ser"}\NormalTok{),}
  \AttributeTok{report =} \StringTok{"vc*"}\NormalTok{,}
  \AttributeTok{digits =} \DecValTok{3}\NormalTok{,}
  \AttributeTok{no.space =} \ConstantTok{TRUE}\NormalTok{,}
  \AttributeTok{align =} \ConstantTok{TRUE}\NormalTok{,}
  \AttributeTok{notes =} \StringTok{"Standard errors in parentheses. *** p\textless{}0.01, ** p\textless{}0.05, * p\textless{}0.1"}
\NormalTok{)}
\end{Highlighting}
\end{Shaded}

\begin{verbatim}
## 
## <table style="text-align:center"><caption><strong>Immigrant-Only Models: Retirement Effects</strong></caption>
## <tr><td colspan="4" style="border-bottom: 1px solid black"></td></tr><tr><td style="text-align:left"></td><td colspan="3"><em>Dependent variable:</em></td></tr>
## <tr><td></td><td colspan="3" style="border-bottom: 1px solid black"></td></tr>
## <tr><td style="text-align:left"></td><td colspan="3">Retired</td></tr>
## <tr><td style="text-align:left"></td><td colspan="2"><em>OLS</em></td><td><em>logistic</em></td></tr>
## <tr><td style="text-align:left"></td><td>Baseline</td><td>DID (LPM)</td><td>DID (Logit)</td></tr>
## <tr><td style="text-align:left"></td><td>(1)</td><td>(2)</td><td>(3)</td></tr>
## <tr><td colspan="4" style="border-bottom: 1px solid black"></td></tr><tr><td style="text-align:left">Post</td><td>0.583<sup>***</sup></td><td>0.564<sup>***</sup></td><td>2.590<sup>***</sup></td></tr>
## <tr><td style="text-align:left">Eligible</td><td>0.013<sup>***</sup></td><td>0.009<sup>**</sup></td><td>0.062<sup>**</sup></td></tr>
## <tr><td style="text-align:left">Post × Eligible</td><td></td><td>0.020<sup>*</sup></td><td>0.098</td></tr>
## <tr><td style="text-align:left">Constant</td><td>0.174<sup>***</sup></td><td>0.177<sup>***</sup></td><td>-1.538<sup>***</sup></td></tr>
## <tr><td colspan="4" style="border-bottom: 1px solid black"></td></tr><tr><td style="text-align:left">Observations</td><td>105,133</td><td>105,133</td><td>105,133</td></tr>
## <tr><td style="text-align:left">R<sup>2</sup></td><td>0.316</td><td>0.316</td><td></td></tr>
## <tr><td style="text-align:left">Adjusted R<sup>2</sup></td><td>0.316</td><td>0.316</td><td></td></tr>
## <tr><td style="text-align:left">Log Likelihood</td><td></td><td></td><td>-52,369.430</td></tr>
## <tr><td style="text-align:left">Akaike Inf. Crit.</td><td></td><td></td><td>104,746.900</td></tr>
## <tr><td colspan="4" style="border-bottom: 1px solid black"></td></tr><tr><td style="text-align:left"><em>Note:</em></td><td colspan="3" style="text-align:right"><sup>*</sup>p<0.1; <sup>**</sup>p<0.05; <sup>***</sup>p<0.01</td></tr>
## <tr><td style="text-align:left"></td><td colspan="3" style="text-align:right">Standard errors in parentheses. *** p<0.01, ** p<0.05, * p<0.1</td></tr>
## </table>
\end{verbatim}

Below is a well‐structured project outline incorporating your ideas.
I've also added a few suggestions to strengthen your analysis.

\begin{center}\rule{0.5\linewidth}{0.5pt}\end{center}

\section{Project Outline: The Impact of OAS Eligibility on
Retirement}\label{project-outline-the-impact-of-oas-eligibility-on-retirement}

\subsection{I. Introduction}\label{i.-introduction}

\begin{itemize}
\tightlist
\item
  \textbf{Overview:}

  \begin{itemize}
  \tightlist
  \item
    Briefly describe Canada's three‐pillar retirement system with
    emphasis on OAS (the foundational benefit) and CPP (the work-related
    benefit).
  \item
    State the motivation: to estimate the causal impact of OAS
    eligibility on retirement behavior (measured as labour force
    participation), while accounting for ``pollution'' from CPP and
    other income sources.
  \end{itemize}
\item
  \textbf{Research Focus:}

  \begin{itemize}
  \tightlist
  \item
    Primary focus on OAS eligibility using the age-65 threshold.
  \item
    Secondary analysis on how contamination from CPP and differential
    immigrant eligibility might affect the results.
  \end{itemize}
\end{itemize}

\subsection{II. Research Questions}\label{ii.-research-questions}

\begin{itemize}
\tightlist
\item
  \textbf{Main Question:}

  \begin{itemize}
  \tightlist
  \item
    ``What is the causal effect of OAS eligibility (triggered at age 65)
    on retirement behavior among Canadians aged 55 to 74?''
  \end{itemize}
\item
  \textbf{Subgroup/Extension Questions:}

  \begin{itemize}
  \tightlist
  \item
    ``How does the effect of OAS eligibility differ between domestic
    Canadians and immigrants, especially when immigrants' eligibility
    depends on their age at arrival?''
  \item
    ``How correlated are CPP receipt and OAS eligibility, and how might
    CPP or other income sources influence observed retirement
    patterns?''
  \end{itemize}
\end{itemize}

\subsection{III. Data and Sample
Construction}\label{iii.-data-and-sample-construction}

\begin{itemize}
\tightlist
\item
  \textbf{Data Source:}

  \begin{itemize}
  \tightlist
  \item
    2021 Canadian Census.
  \end{itemize}
\item
  \textbf{Sample Restrictions:}

  \begin{itemize}
  \tightlist
  \item
    \textbf{Age Range:} Focus on individuals aged 55--74 (corresponding
    to age groups: 15: 55--59, 16: 60--64, 17: 65--69, 18: 70--74).
  \item
    \textbf{Labour Force Information:} Exclude those who have never
    worked or did not report a labour force status.
  \item
    \textbf{Immigration Status:}

    \begin{itemize}
    \tightlist
    \item
      Exclude non-permanent residents or those with missing immigration
      status.
    \item
      Construct subgroups based on immigration age:

      \begin{itemize}
      \tightlist
      \item
        \textbf{Group A:} Immigrants who arrived before age 55 (more
        likely to be eligible for OAS).
      \item
        \textbf{Group B:} Recent immigrants (who are generally
        ineligible for OAS due to insufficient residency).
      \item
        \textbf{SPECIFICALLY:} the immigrants can be separated into 2
        groups: e.g., when looking at age group 55-59, there are two
        ways to set the control group.
      \end{itemize}

      \begin{enumerate}
      \def\labelenumi{\arabic{enumi}.}
      \tightlist
      \item
        Strict: Set those who within the AGEGRP of 55-59 to immigrate
        only in immage 50-54, so the longest year they would stay in
        Canada is 9 years, but less sample may be available;
      \item
        Fuzzy: Set IMMAGE to 45-54, where the longest year they would
        stay in Canada is 14 years, but will have relatively more sample
        included. We are going to run our model on both of them.
      \end{enumerate}
    \end{itemize}
  \end{itemize}
\end{itemize}

\subsection{IV. Treatment, Control, and Identification
Strategy}\label{iv.-treatment-control-and-identification-strategy}

\begin{itemize}
\tightlist
\item
  \textbf{Treatment Definition:}

  \begin{itemize}
  \tightlist
  \item
    \textbf{Primary Treatment:} Eligibility for OAS, which is assigned
    at age 65 for individuals who have lived in Canada for at least 10
    years.
  \item
    \textbf{Verification Steps:}

    \begin{itemize}
    \tightlist
    \item
      Test whether there is a clear cutoff---confirm that no one below
      age 65 receives OAS.
    \item
      Check if all Canadian-born individuals above 65 are indeed
      eligible.
    \end{itemize}
  \end{itemize}
\item
  \textbf{Control Groups for DID Analysis:}

  \begin{itemize}
  \tightlist
  \item
    \textbf{Within Age Groups:}

    \begin{itemize}
    \tightlist
    \item
      Use pre-eligibility age groups (55--59 and 60--64; i.e., AGEGRP 15
      \& 16) to establish a baseline common trend.
    \item
      Compare with post-eligibility groups (65--69 and 70--74; i.e.,
      AGEGRP 17 \& 18).
    \end{itemize}
  \item
    \textbf{By Immigration Status:}

    \begin{itemize}
    \tightlist
    \item
      Compare domestic Canadians (and immigrants who arrived early) as
      the treatment group to recent immigrants (those immigrating later,
      with AGEIMM group ≥ 12) who are not eligible.
    \end{itemize}
  \end{itemize}
\item
  \textbf{Identification Strategy:}

  \begin{itemize}
  \tightlist
  \item
    Employ a Difference-in-Differences (DID) framework:

    \begin{itemize}
    \tightlist
    \item
      \textbf{Assumption:} The common trend assumption holds over a
      narrow age range. Here, the pre-treatment trend (observed in age
      groups 15 and 16) provides the counterfactual for what would have
      happened in the absence of OAS eligibility.
    \item
      \textbf{Addressing ``Pollution'':} Include covariates such as CPP
      receipt and other demographic factors to mitigate potential
      confounding from alternative income sources.
    \end{itemize}
  \end{itemize}
\end{itemize}

\subsection{V. Common Trends
Assumption}\label{v.-common-trends-assumption}

\begin{itemize}
\tightlist
\item
  \textbf{Testing Common Trends:}

  \begin{itemize}
  \tightlist
  \item
    Use age groups 15 (55--59) and 16 (60--64) to check if pre-treatment
    trends in labor force participation are similar. The common trend
    will simply be the slope of the line connecting AGEGRP of
    55\textasciitilde59 \& 60\textasciitilde64.
  \item
    \textbf{IMPORTANT} Here we actually can extend our CTA by looking at
    other early immigrants who are not eligible (i.e., not have stayed
    in Canada for 10 years, for those who are currently 45-54 and
    immigrated within 10 years) to more generally test the CTA, but we
    have to note here that these people who are not eligible for OAS now
    will be eligible for OAS in the future, which may embed some bias in
    the CTA of our current analysis. This may make us need to adjust the
    previous AGEGRP that we were focusing on a bit.
  \end{itemize}
\item
  \textbf{Discussion:}

  \begin{itemize}
  \tightlist
  \item
    Explain how the narrow age range minimizes potential differences in
    unobserved characteristics.
  \item
    Discuss limitations of a single-year (2021) cross-sectional dataset
    and how careful subgroup selection (especially among immigrants) can
    bolster the assumption.
  \end{itemize}
\end{itemize}

\subsection{VI. Empirical Models and Analysis
Plan}\label{vi.-empirical-models-and-analysis-plan}

\begin{itemize}
\tightlist
\item
  \textbf{Model 1 -- Main DID Analysis:}

  \begin{itemize}
  \tightlist
  \item
    \textbf{Specification:}

    \begin{itemize}
    \tightlist
    \item
      Compare labour force participation before and after the age-65
      cutoff.
    \item
      Model includes age-group dummies, treatment indicator (OAS
      eligibility), and interaction terms.
    \end{itemize}
  \item
    \textbf{Key Parameter:}

    \begin{itemize}
    \tightlist
    \item
      The DID coefficient estimating the jump in retirement (or drop in
      labour force participation) at age 65.
    \end{itemize}
  \end{itemize}
\item
  \textbf{Model 2 -- Subgroup Analysis by Immigration Status:}

  \begin{itemize}
  \tightlist
  \item
    \textbf{Purpose:}

    \begin{itemize}
    \tightlist
    \item
      To compare outcomes between domestic Canadians (and immigrants who
      arrived early) and recent immigrants.
    \end{itemize}
  \item
    \textbf{Specification:}

    \begin{itemize}
    \tightlist
    \item
      Same DID framework but run separately (or include interaction
      terms with immigration status) to test if the OAS effect is
      stronger among those eligible.
    \end{itemize}
  \end{itemize}
\item
  \textbf{Additional Analysis:}

  \begin{itemize}
  \tightlist
  \item
    Investigate the correlation between CPP and OAS receipt.
  \item
    Include additional demographic covariates (gender, education,
    income, etc.) to control for confounding factors.
  \end{itemize}
\end{itemize}

\subsection{VII. Robustness Checks and
Extensions}\label{vii.-robustness-checks-and-extensions}

\begin{itemize}
\tightlist
\item
  \textbf{Robustness Checks:}

  \begin{itemize}
  \tightlist
  \item
    Vary the age range definitions slightly to test sensitivity.
  \item
    Test alternative specifications (e.g., regression analysis with
    additional covariates, or logistic versus linear probability
    models).
  \end{itemize}
\item
  \textbf{Extensions:}

  \begin{itemize}
  \tightlist
  \item
    Explore an instrumental variable approach: find an instrument that
    isolates the exogenous variation in OAS receipt. Specifically, here
    we are looking at a more specific group of people: the compliers,
    who will change their behavior due to the policy. This can be
    relatively easy to calculate if we use the result of DD (the
    \textbf{intend to treatment}) and divide it by the first stage.
  \item
    Analyze additional outcomes related to retirement quality or
    financial well-being.
  \end{itemize}
\end{itemize}

\subsection{VIII. Expected Contributions and
Limitations}\label{viii.-expected-contributions-and-limitations}

\begin{itemize}
\tightlist
\item
  \textbf{Contributions:}

  \begin{itemize}
  \tightlist
  \item
    Provide causal evidence on the impact of OAS eligibility on
    retirement decisions using a quasi-experimental design.
  \item
    Offer insights into how immigrant status and the timing of
    immigration affect eligibility and retirement behavior.
  \end{itemize}
\item
  \textbf{Limitations:}

  \begin{itemize}
  \tightlist
  \item
    Use of cross-sectional data from a single year limits the ability to
    observe long-term trends.
  \item
    Potential residual confounding from other pension systems (CPP,
    private savings) that must be carefully addressed in the analysis.
  \end{itemize}
\end{itemize}

\subsection{IX. Conclusion}\label{ix.-conclusion}

\begin{itemize}
\tightlist
\item
  \textbf{Summary:}

  \begin{itemize}
  \tightlist
  \item
    Recap the project's objectives, the identification strategy using
    DID, and the role of different subgroups.
  \end{itemize}
\item
  \textbf{Policy Implications:}

  \begin{itemize}
  \tightlist
  \item
    Discuss how findings might inform public policy regarding retirement
    benefits and pension reforms in Canada.
  \end{itemize}
\item
  \textbf{Next Steps:}

  \begin{itemize}
  \tightlist
  \item
    Outline further research avenues or potential data improvements for
    future studies.
  \end{itemize}
\end{itemize}

\begin{center}\rule{0.5\linewidth}{0.5pt}\end{center}

\subsubsection{Final Comments}\label{final-comments}

\begin{itemize}
\tightlist
\item
  \textbf{Detail the Data Cleaning Process:} Clearly outline how you
  will handle missing data and verify the OAS cutoff.
\item
  \textbf{Clarify Variable Construction:} Explicitly define your outcome
  (e.g., ``Not in Labour Force'' indicator), treatment, and covariates.
\item
  \textbf{Consider Graphical Evidence:} In addition to DID estimates,
  visualizing trends across age groups (e.g., with 100\% stacked bar
  charts or RD plots) can bolster your common trends assumption.
\end{itemize}

This outline should provide a strong roadmap for your project while
remaining flexible enough to incorporate additional insights as you
proceed.

\begin{center}\rule{0.5\linewidth}{0.5pt}\end{center}

\end{document}
